\section*{Chapitre \ref{ch:b2:alkene-redox} --- Oxydation et réduction des alcènes}

\begin{solution}{exo:b2:alkene-redox:1}
Méthylcyclohexane~; 1-méthylcyclohexan-1-ol~; trans-2-méthylcyclohexan-1-ol (addition syn de
H et de B, donc le \ce{OH} et le nouvel H du même côté, le méthyle en trans du \ce{OH})~;
oxyde de 1-méthylcyclohexène (\omterm{def:b2:alkene-redox:epoxide}{époxyde})~; cis-1-méthylcyclohexane-1,2-diol~; 6-oxoheptanal
\ce{CH3CO(CH2)4CHO}.
\end{solution}

\begin{solution}{exo:b2:alkene-redox:2}
Hydroboration--oxydation~: 2-méthylpropan-1-ol. Hydratation en milieu acide~:
2-méthylpropan-2-ol.
\end{solution}

\begin{solution}{exo:b2:alkene-redox:3}
trans-2,3-Diméthyloxirane, les deux méthyles de part et d'autre du cycle. Il est chiral (pas
de plan miroir) et se forme sous forme de mélange racémique.
\end{solution}

\begin{solution}{exo:b2:alkene-redox:4}
Traitement réducteur~: deux molécules de propanal. Traitement oxydant~: deux molécules d'acide
propanoïque.
\end{solution}

\begin{solution}{exo:b2:alkene-redox:5}
(a) méso-butane-2,3-diol, inactif parce qu'il est achiral. (b) les diols
(2\textit{R},3\textit{R}) et (2\textit{S},3\textit{S}) en quantités égales, inactifs
parce que racémiques.
\end{solution}

\begin{solution}{exo:b2:alkene-redox:6}
En milieu acide, l'eau attaque le carbone tertiaire~; en milieu basique, l'hydroxyde attaque le \ce{CH2}~; dans les deux
cas, un \ce{OH} se retrouve sur chaque carbone~: le même 2-méthylpropane-1,2-diol. Avec le méthanol, les deux
voies donnent des éthers différents (comme dans le chapitre).
\end{solution}

\begin{solution}{exo:b2:alkene-redox:7}
Hydrogénation par un équivalent de \ce{H2} sur un catalyseur au palladium empoisonné
(Lindlar)~: l'addition syn sur la triple liaison s'arrête à l'alcène (\textit{Z}).
\end{solution}

\begin{solution}{exo:b2:alkene-redox:8}
La propanone (3 C) et le butanal (4 C) se relient par leurs carbones carbonyles~: c'est le
2-méthylhex-2-ène, \ce{(CH3)2C=CH-CH2CH2CH3}.
\end{solution}

\begin{solution}{exo:b2:alkene-redox:9}
\ce{BH3} (ou un borane encombré), puis \ce{H2O2} et \ce{NaOH}~: hexan-1-ol. L'hydratation en milieu acide
passe par le carbocation secondaire et donne l'hexan-2-ol (Markovnikov), avec des
sous-produits de réarrangement.
\end{solution}

\begin{solution}{exo:b2:alkene-redox:10}
Le \omterm{def:b2:alkene-redox:epoxide}{peroxyacide}, électrophile, réagit plus vite avec la double liaison cyclique, plus substituée et plus riche
en électrons. Le borane encombré réagit avec le \ce{C=CH2} terminal, moins encombré.
\end{solution}

\begin{solution}{exo:b2:alkene-redox:11}
trans~: \omterm{def:b2:alkene-redox:epoxide}{peroxyacide}, puis eau en milieu acide (ouverture anti). cis~: \ce{OsO4} catalytique avec un
réoxydant, ou permanganate dilué à froid (syn).
\end{solution}

\begin{solution}{exo:b2:alkene-redox:12}
\ce{C6H10} possède deux insaturations~; un équivalent de \ce{H2} signifie une \ce{C=C}, donc
un cycle. Un seul dialdéhyde à six carbones provient d'un cycle ouvert au niveau de sa double liaison~:
le cyclohexène.
\end{solution}

\begin{solution}{pb:b2:alkene-redox:1}
\textbf{1.} $10 \times 12.011 + 16 \times 1.008 = \qty{136.24}{g/mol}$.
\textbf{2.} $(2 \times 10 + 2 - 16)/2 = 3$.
\textbf{3.} Un cycle et deux doubles liaisons.
\textbf{4.} \ce{C10H16 + 2H2 -> C10H20}, le 1-isopropyl-4-méthylcyclohexane.
\textbf{5.} $1.00/136.24 = \qty{7.34}{mmol}$ de limonène, \qty{14.68}{mmol} de \ce{H2}.
\textbf{6.} $V = nRT/p = 0.01468 \times 8.314 \times 298.15/10^5 = \qty{3.64e-4}{m^3}$,
\qty{364}{mL}.
\textbf{7.} Non~: les carbones 1 et 4 du cycle portent chacun deux chemins de cycle identiques~; les isomères cis et
trans sont achiraux.
\textbf{8.} Le méthanal \ce{HCHO} (1 C), issu de l'extrémité \ce{=CH2}, et un seul composé en \ce{C9},
le 3-acétyl-6-oxoheptanal \ce{CH3CO-CH2CH2-CH(COCH3)-CH2-CHO}.
\textbf{9.} Ses deux carbones appartiennent au même cycle~: rompre la double liaison ouvre le cycle
sans rien séparer.
\textbf{10.} Le méthanal (un gaz).
\textbf{11.} Les groupes aldéhyde deviennent des acides carboxyliques (le méthanal va jusqu'à \ce{CO2})~; les
cétones restent inchangées.
\textbf{12.} La cétone et l'aldéhyde de la chaîne en \ce{C9} étaient reliés dans le cycle~; la seconde
cétone (\ce{COCH3}) et le méthanal étaient les deux extrémités de la double liaison de la chaîne latérale.
\textbf{13.} Un seul (le C4 du cycle)~; l'\omterm{def:b2:alkene-redox:cleavage}{ozonolyse} ne le touche pas, et il subsiste comme
stéréocentre du produit en \ce{C9}.
\textbf{14.} Le \ce{C=CH2} exocyclique, moins encombré que la double liaison cyclique trisubstituée.
\textbf{15.} Le 2-(4-méthylcyclohex-3-én-1-yl)propan-1-ol~: le \ce{CH2OH} sur l'ancien carbone
terminal.
\textbf{16.} Primaire.
\textbf{17.} Le \ce{OH} se place sur le carbone le moins substitué, à l'opposé de l'orientation
prévue par la règle de Markovnikov pour l'hydratation.
\textbf{18.} Deux configurations au nouveau centre, ce qui donne deux diastéréoisomères (avec le centre
déjà présent).
\textbf{19.} L'hydratation en milieu acide (ou une autre hydratation de Markovnikov)~: l'alcool tertiaire,
l'$\alpha$-terpinéol.
\textbf{20.} La double liaison cyclique~: trisubstituée, plus riche en électrons.
\textbf{21.} Le 1,2-époxy-1-méthyl-4-(prop-1-én-2-yl)cyclohexane, l'oxygène pontant C1 et C2.
\textbf{22.} Sur C1, le carbone le plus substitué (celui qui porte le méthyle).
\textbf{23.} Le 1-méthyl-4-(prop-1-én-2-yl)cyclohexane-1,2-diol, les deux \ce{OH} en trans (ouverture
anti).
\textbf{24.} Une \omterm{def:b2:alkene-redox:dihydroxylation}{dihydroxylation} syn de la double liaison cyclique (\ce{OsO4} catalytique), qui donne le diol
cis --- à condition de pouvoir faire réagir le réactif sélectivement avec cette double liaison.
\textbf{25.} \qty{1.00}{g} de limonène absorbe $\boldsymbol{\approx \qty{364}{mL}}$ de
dihydrogène.
\end{solution}
